\documentclass{situationist-preprint}

\title{Machine-Generated Art and the Institutional Theory}
\msid{AES-15-a}
\msfamily{AES-15}
\msdiscipline{Aesthetics}
\mstype{Principal manuscript}
\mskeywords{machine art; institutional theory; artworld; uptake; art-attempts; authorship}

\begin{document}
\maketitle

\begin{abstract}
Can machines make art? We argue that the answer, on the account of art defended in AES-01, is a qualified yes: machine outputs can be art if they are presented as art within an artworld and taken up by its critical practice, and whether they are so taken up is an open question that philosophy cannot settle. The argument separates three questions usually run together: whether the output of a machine can be an artwork, whether the machine is its artist, and whether the output is good art. On the institutional account, the first question turns on presentation and uptake, neither of which requires that the system producing the object have intentions. We reply to the objection that art requires an art-attempt by its maker, to the objection that the presenter rather than the machine is the artist, and to the objection that market prices already show uptake. The result is deliberately modest: it establishes a possibility and locates the remaining question in the artworld rather than in the metaphysics of minds.
\end{abstract}

\section{The question}

Image generators, text models and music systems now produce outputs that are exhibited, sold and reviewed. The question whether such outputs can be art is often answered by appeal to the question whether machines can have intentions, feelings or creativity \citep{Boden2010,Coeckelbergh2017,Hertzmann2018}. We argue that, on the institutional account of art, this route is a detour. The status of an object as art does not depend on the inner states of whatever produced it, but on what is done with it.

We distinguish three questions:
\begin{enumerate}
\item[(Q1)] Can a machine output be an artwork?
\item[(Q2)] Is the machine the artist of that work?
\item[(Q3)] Is the work good?
\end{enumerate}
This manuscript answers (Q1). It leaves (Q3) to criticism and says only enough about (Q2) to keep it apart from (Q1).

\section{Framework}

We take over the account of AES-01, which descends from \citet{Danto1964} and \citet{Dickie1974}.

\begin{definition}\label{d:presented}
An object is \emph{presented as art} within an artworld if some participant in that artworld offers it to a public as a candidate for the kind of appreciation and criticism the artworld practices.
\end{definition}

\begin{definition}\label{d:uptake}
A presented object receives \emph{critical uptake} if the artworld's critical practice engages it as a candidate: it is interpreted, compared with other works, and assessed by the standards the practice applies to artworks.
\end{definition}

\begin{principle}[Institutional condition]\label{p:inst}
An object is an artwork if and only if it is presented as art within an artworld and receives critical uptake there (AES-01).
\end{principle}

Two features of \ref{p:inst} matter here. First, the presenter need not be the maker: found objects and readymades are presented by someone who did not fabricate them. Second, uptake is not mere attention. A price, a headline or a viral circulation is not yet critical engagement.

\section{The argument}

\begin{proposition}\label{prop:possible}
Nothing in \ref{p:inst} excludes machine outputs. Hence a machine output can be an artwork.
\end{proposition}

\begin{proof}[Argument]
The two conditions of \ref{p:inst} concern the relation between an object and an artworld. Presentation (Definition~\ref{d:presented}) requires a participant in the artworld who offers the object; it places no condition on how the object came to exist. A driftwood sculpture, a readymade and a chance composition were each produced by processes without artistic intentions, and each has been presented. Uptake (Definition~\ref{d:uptake}) is a condition on the artworld's practice, not on the object's origin. So an object produced by a machine can satisfy both conditions, provided someone presents it and the practice engages it. Since the conditions are jointly sufficient, the output can be an artwork.
\end{proof}

\begin{proposition}\label{prop:open}
Whether machine outputs are in fact artworks is not settled by Proposition~\ref{prop:possible}, and is at present open.
\end{proposition}

\begin{proof}[Argument]
\ref{p:inst} makes art status depend on uptake, and uptake is a matter of what a critical practice does. Proposition~\ref{prop:possible} shows only that nothing in the account forbids it. Whether the practice engages machine outputs as candidates---interpreting them, situating them in traditions, judging them by the standards it applies to other works---is a fact about the practice. The evidence so far is mixed: such outputs are exhibited and sold, and some critics review them as works, while others treat them as products of a technique whose interest lies in the technique. A question whose answer depends on how a practice will settle is not answered by an analysis of the concept the practice uses.
\end{proof}

The two propositions together give the conclusion of AES-15: machine outputs can be art if they are taken up by an artworld, and whether they are so taken up remains open.

\section{Objections and replies}

\begin{objection}
Art requires an art-attempt: something is an artwork only if someone tried to make an artwork \citep{MagUidhir2013}. A machine without intentions cannot attempt anything, so its outputs cannot be art.
\end{objection}

\begin{reply}
The intention that matters on the institutional account is the presenter's, not the producer's. Even if art requires an attempt, the presenter who selects, frames and offers an output to the artworld makes one. The objection, if sound, shows that machine outputs are not art without presentation; it does not show that they cannot become art through it. The same structure applies to readymades, where the fabricator of the urinal attempted nothing artistic.
\end{reply}

\begin{objection}
If presentation does the work, then the artist is the presenter, and the machine contributes nothing more than a camera or a brush. The question was whether machines make art, and the answer turns out to be that people do.
\end{objection}

\begin{reply}
This is (Q2), not (Q1). Proposition~\ref{prop:possible} concerns whether the output can be an artwork, which it can whoever the artist turns out to be. On (Q2) we note only that the presenter's contribution to a machine output differs from the photographer's to a photograph: the presenter often does not fix the composition, which may be the product of a process no one controls in detail. Whether that difference makes the machine a co-author is a question about authorship, connected to the question of agency treated in ACT-16, on which current systems are not responsible agents and responsibility for their outputs lies with those who deploy them. Authorship and art status can come apart.
\end{reply}

\begin{objection}
Uptake has already occurred. A portrait produced by a generative network was sold at auction by Christie's in 2018 for a sum far above its estimate.
\end{objection}

\begin{reply}
A sale is market uptake, not critical uptake (Definition~\ref{d:uptake}). The auction showed that the object could be bought and sold as a collectible; whether it was interpreted and assessed as a work, rather than as an event in the history of a technology, is precisely what remains disputed. The case supports Proposition~\ref{prop:open}: the artworld's response is divided.
\end{reply}

\begin{objection}
The institutional theory is too permissive. If uptake can make anything art, the answer to (Q1) is trivial and uninformative.
\end{objection}

\begin{reply}
The answer is informative because the alternative views make it false. On accounts that require the maker's expressive or intentional states, machine outputs could not be art however they were received. The institutional account relocates the question from the producer's mind to the artworld's practice, and that relocation is the substantive result. It is not trivial that the question cannot be settled by establishing what machines can feel.
\end{reply}

\section{What the result does not establish}

It does not establish that any machine output is art, that machines are artists, or that machine outputs are good. It does not depend on whether machines can think, a question taken up in MND-08; if they can, a machine might also be a presenter, but that possibility is not needed here.

\section{Conclusion}

On the institutional account, the question whether machines can make art is a question about artworlds, not about machines. Machine outputs can be art; whether they are depends on whether the artworld takes them up as it takes up other works, and that is not yet decided.

\begin{thebibliography}{9}
\bibitem[Boden(2010)]{Boden2010} Boden, M.~A. (2010). \emph{Creativity and Art: Three Roads to Surprise}. Oxford University Press.
\bibitem[Coeckelbergh(2017)]{Coeckelbergh2017} Coeckelbergh, M. (2017). Can machines create art? \emph{Philosophy \& Technology}, 30(3), 285--303.
\bibitem[Danto(1964)]{Danto1964} Danto, A.~C. (1964). The artworld. \emph{The Journal of Philosophy}, 61(19), 571--584.
\bibitem[Dickie(1974)]{Dickie1974} Dickie, G. (1974). \emph{Art and the Aesthetic: An Institutional Analysis}. Cornell University Press.
\bibitem[Hertzmann(2018)]{Hertzmann2018} Hertzmann, A. (2018). Can computers create art? \emph{Arts}, 7(2), 18.
\bibitem[Mag Uidhir(2013)]{MagUidhir2013} Mag Uidhir, C. (2013). \emph{Art and Art-Attempts}. Oxford University Press.
\end{thebibliography}

\end{document}
